\documentclass{beamer}
\usepackage[utf8]{inputenc}

\usetheme{Madrid}
\usecolortheme{default}
\usepackage{amsmath,amssymb,amsfonts,amsthm}
\usepackage{txfonts}
\usepackage{tkz-euclide}
\usepackage{listings}
\usepackage{adjustbox}
\usepackage{array}
\usepackage{tabularx}
\usepackage{gvv}
\usepackage{lmodern}
\usepackage{circuitikz}
\usepackage{tikz}
\usepackage{graphicx}
\usepackage{mathtools}
\usepackage{subcaption}
\setbeamertemplate{page number in head/foot}[totalframenumber]

\usepackage{tcolorbox}
\tcbuselibrary{minted,breakable,xparse,skins}



\definecolor{bg}{gray}{0.95}
\DeclareTCBListing{mintedbox}{O{}m!O{}}{%
  breakable=true,
  listing engine=minted,
  listing only,
  minted language=#2,
  minted style=default,
  minted options={%
    linenos,
    gobble=0,
    breaklines=true,
    breakafter=,,
    fontsize=\small,
    numbersep=8pt,
    #1},
  boxsep=0pt,
  left skip=0pt,
  right skip=0pt,
  left=25pt,
  right=0pt,
  top=3pt,
  bottom=3pt,
  arc=5pt,
  leftrule=0pt,
  rightrule=0pt,
  bottomrule=2pt,
  toprule=2pt,
  colback=bg,
  colframe=orange!70,
  enhanced,
  overlay={%
    \begin{tcbclipinterior}
    \fill[orange!20!white] (frame.south west) rectangle ([xshift=20pt]frame.north west);
    \end{tcbclipinterior}},
  #3,
}
\lstset{
    language=C,
    basicstyle=\ttfamily\small,
    keywordstyle=\color{blue},
    stringstyle=\color{orange},
    commentstyle=\color{green!60!black},
    numbers=left,
    numberstyle=\tiny\color{gray},
    breaklines=true,
    showstringspaces=false
}

\title 
{1.2.40}
\date{\today}


\author 
{Vivek K Kumar - EE25BTECH11062}



\begin{document}


\frame{\titlepage}

\begin{frame}{Question}
\textbf{Question}:\\
\begin{align}
    \lim\limits_{x\to0} \frac{e^x-\brak{1+x+\frac{x^2}{2}}}{x^3} = 
\end{align}
\begin{enumerate}
    \item $0$
    \item $1/6$
    \item $1/3$
    \item $1$
\end{enumerate}
\end{frame}

\begin{frame}[fragile, allowframebreaks]{Solution}
\textbf{Method 1 - Taylor series}\\
Taylor series expansion of $e^x$
\begin{align}
    e^x = \sum\limits_{n=0}^\infty \frac{x^n}{n!} = 1+x+\frac{x^2}{2}+\frac{x^3}{6}+\frac{x^4}{24} \dots
\end{align}

Then the limit becomes
\begin{align}
    &\lim\limits_{x\to0}\frac{\brak{1+x+\frac{x^2}{2}+\frac{x^3}{6}+\frac{x^4}{24} \dots} - \brak{1+x+\frac{x^2}{2}}}{x^3} \\
    &= \lim\limits_{x\to0} \frac{\frac{x^3}{6} + \frac{x^4}{24} + \dots}{x^3} \\
    &= \lim\limits_{x\to0} \frac{1}{6} + \frac{x}{24} + \frac{x^2}{120}\dots \\
   & = \boxed{\frac{1}{6}}
\end{align}
\begin{center}
$\boxed{\mathbf{ \lim\limits_{x\to0} \frac{e^x-\brak{1+x+\frac{x^2}{2}}}{x^3} = \frac{1}{6}}}$
\end{center}

\textbf{Method 2 - L'Hospital Rule}
Since numerator and denominator tends to 0, we can apply L'Hospital rule. 

\begin{align}
    &\lim\limits_{x\to0} \frac{e^x-\brak{1+x+\frac{x^2}{2}}}{x^3} \\
    &= \lim\limits_{x\to0}\frac{\frac{d}{dx}\brak{e^x - \brak{1+x+\frac{x^2}{2}}}}{\frac{d}{dx}(x^3)} \\
    &= \lim\limits_{x\to0} \frac{e^x - \brak{1+x}}{3x^2}    
\end{align}

Applying L'hospital again 

\begin{align}
    & \lim\limits_{x\to0} \frac{e^x - \brak{1+x}}{3x^2} \\
    &= \lim\limits_{x\to0}\frac{\frac{d}{dx}\brak{e^x - \brak{1+x}}}{\frac{d}{dx}(3x^2)} \\
    &= \lim\limits_{x\to0} \frac{e^x - 1}{6x} \\
\end{align}

Applying L'hospital one last time
\begin{align}
    & \lim\limits_{x\to0} \frac{e^x - 1}{6x} \\
    &= \lim\limits_{x\to0}\frac{\frac{d}{dx}\brak{e^x - 1}}{\frac{d}{dx}(6x)} \\
    &= \lim\limits_{x\to0} \frac{e^x}{6} \\
    & = \frac{e^0}{6} = \boxed{\frac{1}{6}}
\end{align}

\begin{center}
$\boxed{\textbf{Option 2 is the correct answer}}$
\end{center}
\end{frame}

\begin{frame}[fragile,allowframebreaks]{Python Code - Plotting function }
\begin{lstlisting}[language=Python]
import numpy as np
import matplotlib.pyplot as plt
num = 100
x = np.linspace(-0.5, 0.5, num = num)
y1 = ( np.exp(x) - 1 - x - ((x**2)/2) ) / (x**3)
y2 = 1/6*np.ones(num)
x2 = np.zeros(num)
y3= np.linspace(0.15, 0.20, num=num)
plt.plot(x, y1, label = 'e^x - (1+x+x^2/2) / x^3')
plt.plot(x, y2, label = 'lim x->0 = 1/6 (theoretical)', linestyle = 'dotted')
plt.plot(x2, y3, linestyle='dotted')
plt.legend()
plt.title('Limit visualisation')
\end{lstlisting}
\end{frame}

\begin{frame}{Plot}
    \begin{figure}[H]
    \centering
    \includegraphics[width=0.7\linewidth]{beamer/figs/lim.png}
    \caption{Plot of function}
    \label{fig:sp}
\end{figure}
\end{frame}

\begin{frame}[fragile,allowframebreaks]{Python Code - Visualising L'hospital rule }
\begin{lstlisting}[language=Python]
import numpy as np
import matplotlib.pyplot as plt

num = 100
x = np.linspace(-0.5, 0.5, num = num)

y1 = ( np.exp(x) - 1 - x - ((x**2)/2) ) / (x**3)
y2 = 1/6*np.ones(num)
x2 = np.zeros(num)
y3= np.linspace(0.15, 0.20, num=num)
y4 = ( np.exp(x) - 1 - x ) / (3*x*x)
y5 = ( np.exp(x) - 1 ) / (6*x)
y6 = y2

plt.plot(x, y1, label = 'e^x - (1+x+x^2/2) / x^3')
plt.plot(x, y2, label = 'lim x->0 = 1/6 (theoretical)', linestyle = 'dotted')
plt.plot(x2, y3, linestyle='dotted')
plt.plot(x, y4, label = 'e^x - (1+x) / 3x^2')
plt.plot(x, y5, label = 'e^x - 1 / 6x')
plt.plot(x, y6, label = '1/6')
plt.legend()
plt.title('Lhospital visualisation')
\end{lstlisting}
\end{frame}

\begin{frame}{Plot -2}
    \begin{figure}[H]
    \centering
    \includegraphics[width=0.7\linewidth]{beamer/figs/lhospital.png}
    \caption{L'hospital visualisation}
    \label{fig:spl}
\end{figure}
\end{frame}

\end{document}